\section{Discussion and Limitations}
\label{sec:discussion}

This section reads the challenge's three stated difficulties against our
measurements, states plainly what the explanation output is and is not,
and lists the threats to validity.

\subsection{The three difficulties, measured}
\label{sec:discussion:difficulties}

\textbf{Root cause beyond symptoms.} The challenge describes this as the
hardest part, and our data agrees in a form that can be quantified: in
53\% of real reports, no faulty file is named anywhere in the text, and
on those instances every grep-style method scores zero. Incident memory
is currently the only signal in our system that routinely recovers such
files, because a past fix can point at code that no present log line
mentions. The honest counterpart from \cref{sec:results:negatives} also
stands: on one of four projects, fusing signals still performs worse on
this stratum than content retrieval alone. We also do not attempt
reasoning over control flow, data flow, or dependency structure; memory
is a statistical stand-in for that reasoning, and the challenge's full
causal-analysis ask remains open. The hard stratum is where
future work on this challenge should be judged.

\textbf{Noisy and incomplete log evidence.} Here our benchmark is kinder
than industrial reality, and we say so. Its logs are the excerpts that
reporters chose to paste into an issue, so a human has already filtered
the noise once. A production log bundle with retries, background
activity, and downstream errors is a harder input. The benchmark does
capture the incompleteness side: in many instances the pasted evidence is
a single generic error line, and 90\% of MongoDB reports contain no stack
trace at all. The noise experiment of \cref{sec:results:noise} simulates
the bundle side and bounds it: volume noise is nearly free, and the real
requirement a raw bundle adds is separating interleaved incidents before
matching. Evaluating on true unfiltered runtime
bundles is the clearest next step, and our architecture makes no
assumption that prevents it: every signal consumes plain text.

\textbf{Evidence-to-code semantic retrieval.} Our lexical instantiation
is deliberately simple: BM25 over paths, over file contents within a
per-instance pool, and over past incident reports. It is cheap,
deterministic, and leak-free, and it was enough to beat the strongest
baseline everywhere. Nothing in the fusion design depends on the signals
being lexical. A dense retriever trained for log-to-code matching, or an
LLM reranker over the fused top candidates, drops into the same slot and
inherits the same firewall; we consider this the most promising upgrade
path.

\subsection{What the explanation output is}
\label{sec:discussion:explanation}

The challenge asks for an explanation linking symptoms, log events, and
candidate regions. Our explanation is provenance, not generated text. For
each ranked file the system reports which signals ranked it, the explicit
mentions or matching terms behind the log-to-code signals, and for the
memory signal the keys of the past incidents that voted for it. An
engineer can open the cited prior issues and their fix commits directly.
We prefer this form because every element of it is checkable; a fluent
narrative written by a language model can be added on top, but the
citations are the part an on-call engineer can trust at 3 a.m. Coverage
is complete on the evaluation split: every gold file ranked in the top
five carries at least one checkable element, always including the
matched terms, with an explicit mention for 52\% to 86\% of hits and
past-incident voter keys for 63\% to 80\%. Whether these citations
actually speed engineers up is a user-study question this paper does not
answer.

\subsection{Deployment notes}
\label{sec:discussion:deployment}

Three properties matter for industrial adoption. First, cost: the whole
pipeline is lexical, needs no GPU and no training, and we measured it on
one laptop CPU core with warm caches. The median end-to-end time per
issue, signals plus fusion, is 0.07\,s on Kafka, 0.22\,s on Qt, and
0.41\,s on MariaDB and MongoDB; the 95th percentile stays below 0.7\,s
everywhere, and building the incident-memory index takes at most 0.1\,s
even for the largest pool. Second,
freshness: the memory pool is append-only, and every newly resolved issue
with a linked fix commit widens the searchable history at no retraining
cost. Consistent with this, memory was weakest in the earliest part of
the evaluation period in three of four projects
(\cref{sec:results:rq3}). Third, failure mode: with
an empty or sparse memory pool the system degrades to its log-to-code
signals, which already match the strongest baseline; memory is an
additive signal behind a firewall, never a dependency.

\subsection{Threats to validity}
\label{sec:discussion:threats}

\textbf{Construct.} Mined gold inherits the imperfections of
\cref{sec:benchmark:limits}: a fix commit's files are the repair site,
which is occasionally not the origin site, and commit-message linkage can
be noisy despite our filters. The module metric depends on a per-project
depth choice, which we fixed before evaluation and release with the
artifact.

\textbf{Internal.} The development split precedes the evaluation split in
time, which protects against tuning on the future but, as
\cref{sec:setup:protocol} records, biases tuned weights against memory.
We tested how much this one-shot tuning costs with a rolling-origin
variant that retunes the conditioned weights at every quarter of the
evaluation period on all history available at that point: overall Acc@5
moves by $-0.007$ to $+0.021$ across the four projects. Tuning once on
the earliest 30\% is therefore stable rather than strongly conservative,
and a maintained deployment that retunes quarterly should expect
essentially the reported numbers. The regime-conditioned
tuning is fragile when a regime has few development instances, as Qt
shows.

\textbf{External.} Four projects, two tracker platforms, and one
ecosystem of conventions (issue keys in commit messages) do not cover all
industrial settings. Projects without disciplined issue--commit linkage
cannot use our mining recipe directly, although organizations with any
form of change tracking can substitute their own linkage. Logs in the
benchmark are reporter-curated, as discussed above.
